\chapter{Binary File Format Specification}
\label{ch:spec}

For Mesh-LZ to function as an independent, interoperable image format across disparate hardware architectures, a strictly defined binary specification is required. This chapter formalizes the byte-level structure of a Mesh-LZ (\texttt{.mlz}) file. All multi-byte integers are serialized in Little-Endian byte order to align with the native memory architecture of modern x86\_64 and ARM processors, thereby preventing endian-swapping overhead during ultra-fast decode cycles.

\section{File Signature and Header}
Every valid Mesh-LZ file must begin with a 16-byte fixed header.

\begin{table}[h!]
\centering
\begin{tabular}{|l|l|l|l|}
\hline
\textbf{Offset} & \textbf{Length} & \textbf{Type} & \textbf{Description} \\
\hline
0x00 & 4 bytes & ASCII & Magic Signature: \texttt{MLZ!} (\texttt{0x4D 0x4C 0x5A 0x21}) \\
0x04 & 4 bytes & \texttt{u32} & Image Width in pixels \\
0x08 & 4 bytes & \texttt{u32} & Image Height in pixels \\
0x0C & 1 byte & \texttt{u8} & Color Space Mode (See Section \ref{sec:color_mode}) \\
0x0D & 1 byte & \texttt{u8} & Block Size (Typically 8 or 16) \\
0x0E & 2 bytes & Padding & Reserved for future flags (Must be 0x00 0x00) \\
\hline
\end{tabular}
\caption{Mesh-LZ Global File Header}
\end{table}

\subsection{Color Space Modes}
\label{sec:color_mode}
The byte at offset \texttt{0x0C} dictates how the decoder interprets the payload data:
\begin{itemize}
    \item \texttt{0x00}: Truecolor RGB (Lossless)
    \item \texttt{0x01}: YCoCg-R (Reversible, Lossless)
    \item \texttt{0x02}: Indexed Palette (256 colors via NeuQuant)
    \item \texttt{0x03}: YCoCg-R with 4:2:0 Chroma Subsampling (Lossy)
\end{itemize}

\section{Palette Chunk (Optional)}
If the Color Space Mode is set to \texttt{0x02}, a Palette Chunk immediately follows the Global Header. The chunk spans exactly 768 bytes, representing 256 consecutive RGB triplets.
\begin{verbatim}
[R0, G0, B0, R1, G1, B1, ..., R255, G255, B255]
\end{verbatim}
The subsequent block payloads will contain single-byte 8-bit indices rather than full color vectors.

\section{Block Payload Structure}
The remainder of the file is a sequence of compressed blocks, processed in standard row-major raster order (top-left to bottom-right). Because blocks are encoded using Interleaved rANS, their bitstream lengths are entirely variable. Thus, each block is prefixed by a 32-bit length integer.

\begin{table}[h!]
\centering
\begin{tabular}{|l|l|l|l|}
\hline
\textbf{Offset} & \textbf{Length} & \textbf{Type} & \textbf{Description} \\
\hline
0x00 & 4 bytes & \texttt{u32} & Total byte length of the compressed block ($N$) \\
0x04 & 1 byte & \texttt{u8} & Stencil Selection ID (0=Raster, 1=Column, 2=Hilbert) \\
0x05 & $N-1$ bytes & Bitstream & Interleaved rANS encoded byte stream \\
\hline
\end{tabular}
\caption{Mesh-LZ Compressed Block Chunk}
\end{table}

\subsection{Bitstream Payload Details}
The $N-1$ bytes of the bitstream consist of three distinct components interleaved for performance:
\begin{enumerate}
    \item \textbf{LZ Flags:} A bit-packed array indicating whether the following token is a Literal or a Match.
    \item \textbf{Literals:} The raw residual values encoded via rANS.
    \item \textbf{Matches:} Length and distance pairs, where lengths are heavily skewed towards smaller values and distances are constrained to the current $16 \times 16$ block footprint (maximum distance of 255).
\end{enumerate}

To decode the bitstream, the decoder first reads the final 4 bytes to initialize the rANS state integer. It then traverses the bitstream backwards (as is mathematically required by the LIFO nature of Asymmetric Numeral Systems), updating the state and yielding symbols. The generated array of symbols is then inverse-mapped using the Stencil Selection ID to reconstruct the 2D spatial block.

\section{EOF Marker}
To allow streaming implementations to safely terminate without strictly relying on End-Of-File signals from the host OS, Mesh-LZ files append a 4-byte EOF marker consisting of \texttt{0xFF 0xFF 0xFF 0xFF}. If a block size header reads this value, the decoding loop terminates cleanly.
